\documentclass[10pt,a4paper]{amsart}
\usepackage[margin=19mm]{geometry}
\usepackage{amsmath,amssymb,mathtools}
\usepackage[scr=rsfs,cal=euler]{mathalfa}
\usepackage{xfrac}
\usepackage[unicode,hidelinks,bookmarks=false]{hyperref}
\newcommand{\ie}{i.e.\ }
\newcommand{\cf}{cf.\ }
\newcommand{\resp}{respectively\ }
\newcommand{\Subsection}{\subsection}
\providecommand{\nobreakdash}{\nobreak\hbox{-}}
\setlength{\emergencystretch}{2em}
\multlinegap0pt
\usepackage{bm}
\usepackage{bbm}
\usepackage{dsfont}
\usepackage{xspace}
\usepackage{cancel}
\usepackage{mathtools}
\usepackage{amsthm}
\makeatletter
\@ifundefined{Theorem}{\newtheorem{Theorem}{Theorem}}{}
\@ifundefined{Corollary}{\newtheorem{Corollary}[Theorem]{Corollary}}{}
\@ifundefined{Definition}{\newtheorem{Definition}[Theorem]{Definition}}{}
\makeatother
\title[The Gomory-Hu inequality and trees]{The Gomory-Hu inequality and trees}
\author{Oleksiy Dovgoshey, Olga Rovenska}
\date{Source: arXiv:2604.18400v1 (CC BY 4.0)}
\begin{document}
\maketitle

\noindent\emph{Source and licence.} The Gomory-Hu inequality and trees. Version arXiv:2604.18400v1 (CC BY 4.0), distributed under the Creative Commons Attribution 4.0 International licence, as recorded in the arXiv licence metadata for this exact version and shown on the abstract page https://arxiv.org/abs/2604.18400v1. Full rights notice: licenses/arxiv-2604.18400-CC-BY-4.0.txt.\par\smallskip

\noindent\textit{Editorial scope.} Counted unit: the Theorem whose source label is \texttt{tteo5} in the article (source0.tex lines 620--627, source label \texttt{tteo5}), delivered together with the article's own complete original proof of it (source0.tex lines 629--741, one \texttt{proof} environment of 113 lines). No mathematics is written, repaired or re-derived: statement and proof are the article's own text line for line. Required context retained uncounted: kaske (lines 413--419); kasya (lines 426--441); kjvdde (lines 432--439); fdes1 (lines 463--470); zero (lines 521--526); btoy (lines 564--571). Every definition, display and label of those lines is retained unchanged. Omitted from this package and not redistributed: 1--412, 420--425, 440--462, 471--520, 527--563, 572--619, 628--628, 742--1569. Nothing inside a retained excerpt is omitted or altered: every retained body is delivered exactly as the article's lines are written, so no omission receipt of any kind accompanies this package. Citations: the retained excerpts carry no citation command at all, so no bibliography entry of the article is transcribed and no reference is redistributed. Reference adaptation: none was needed and none was made; every reference inside the delivered unit resolves inside it, so no unresolved reference or citation is produced. Printed slips kept literal: no printed word, symbol, hypothesis or number of the counted statement or of its proof was altered. Layout and class: the article's own class is \texttt{sn-jnl} with packages including graphicx, multirow, amsmath, amssymb, amsfonts, amsthm, mathrsfs, appendix, xcolor, textcomp, manyfoot, booktabs, algorithm, algorithmicx; the wrapper is the lane's pinned \texttt{amsart} editorial wrapper at 10pt on A4, which restates the theorem-like environments the retained text needs and adds the article's own macro definitions that the retained text needs (none), with extra wrapper packages (bm, bbm, dsfont, xspace, cancel, mathtools). The delivered numbering is the unit's own. Assets: no retained span contains a figure environment, a graphics-inclusion command, a PDF-page inclusion, a table or a TikZ picture; no image file is copied into this package, the package declares no asset dependency and no image file is redistributed with it. Licence: version arXiv:2604.18400v1 of this article is distributed under the Creative Commons Attribution 4.0 International licence, as recorded in the arXiv licence metadata for this exact version and shown on the abstract page https://arxiv.org/abs/2604.18400v1. A full rights notice is delivered at licenses/arxiv-2604.18400-CC-BY-4.0.txt. Characters: every retained excerpt is pure ASCII, so the delivered engine, XeLaTeX with its default Latin Modern text font, reports no missing character.
\begin{Definition} Let $G(w)$ be a weighted graph.  A weight $w$ is said to be {\it ultrametrizable (pseudoultrametrizable)} if there exists an ultrametric (pseudoultrametric) $\rho_w : V(G) \times V(G) \to \mathbb{R}^+$ such that 
\begin{equation}
\label{kaske}
    w(\{x,y\})=\rho_w(x,y)
\end{equation}
 for each $\{x,y\} \in E(G)$.
\end{Definition}

\begin{Theorem}\label{kasya} Let $G(w)$ be a  weighted graph. The weight $w$ is pseudoultrametrizable if and only if for any cycle $C \subseteq G$ there exist at least two different edges $e_1, e_2 \in E(C)$ such that
\begin{equation}
    \label{samf}
w(e_1) = w(e_2) = \max_{e \in E(C)} w(e).
\end{equation}
If $G$ is connected, $w$ is  pseudoultrametrizable, then the function $\rho_w : V(G) \times V(G) \to \mathbb{R}^+$ defined for all $x,y \in V(G)$ as 
\begin{equation}
    \label{kjvdde}
\rho_w(x,y) :=
\begin{cases}
0, & \text{if } x = y,\\[6pt]
\inf\limits_{P \in \mathcal{P}_{x,y}} \left( \max\limits_{e \in E(P)} w(e) \right), & \text{if } x \ne y,
\end{cases}
\end{equation}
where $\mathcal{P}_{x,y}$ is the set of all paths joining  $x$ and $y$ in $G$, is a pseudoultrametric on $V(G)$ and equality \eqref{kaske} holds.
\end{Theorem}

\begin{equation}
    \label{fdes1}
d_l(x,y) := 
\begin{cases}
0, & \text{if } x = y,\\
\inf\limits_{P\in \mathcal{P}_{x,y}}\left(\max\limits_{v\in V(P)} l(v)\right), & \text{otherwise},
\end{cases}
\end{equation}

\begin{equation}
    \label{zero}
d_l(x,y)
=
\inf_{P \in \mathcal{P}_{x,y}} \left( \max_{e \in E(P)} w(e) \right)
\end{equation}

\begin{Corollary}\label{btoy} Let $G = G(l)$ be a connected labeled graph with  vertex labeling $l : V(G) \to \mathbb R^+,$ and let $d_l : V(G)\times V(G)\to \mathbb R^+$ be generated by $G(l)$. Then the equality
\begin{equation}
    \label{a1}
d_l(x,y)
= \max \{ l(x), l(y) \}
\end{equation}
holds for each edge $ \{x,y\} \in E(G)$.
\end{Corollary}

\begin{Theorem}\label{tteo5} Let $G = G(l)$ be a connected  labeled graph  and let  $d_l : V(G)\times V(G) \to \mathbb{R}^{+}$
be generated by $G(l)$.
Then $d_l$ is  a pseudoultrametric on $V(G)$. Moreover, if $G$ is locally finite, then $d_l$ is an ultrametric if and only if the inequality
\begin{equation}\label{eq:10}
\max\{l(u),l(v)\} > 0
\end{equation}
holds for every edge $\{u,v\} \in E(G)$.
\end{Theorem}

 \begin{proof}


Let us define a weight
\(
w : E(G) \to \mathbb{R}^+
\)
as
\begin{equation}
    \label{trean1}
w(\{u,v\}) := \max \{ l(u), l(v) \}
\end{equation}
for each $\{u,v\} \in E(G).$

First we claim that the weight $w$ is pseudoultrametrizable.

To prove this claim we will use Theorem \ref{kasya}.

Let us consider an arbitrary cycle $C \subseteq G.$
By Theorem \ref{kasya}, $w$ is pseudoultrametrizable if and only if we can find
different $e_1,$ $e_2 \in E(C)$ such that
\begin{equation}
    \label{uro1}
w(e_1) = w(e_2) = \max_{e \in E(C)} w(e).
\end{equation}




Since $C$ is a finite graph, there exists $v_0 \in V(C)$ such that
\begin{equation}
    \label{s1}
l(v_0) = \max_{u \in V(C)} l(u). 
\end{equation}
It follows from the definition of the cycles that there are two distinct vertices $u_1, u_2 \in V(C)$ such that $\{u_1, v_0\}$ and $\{u_2, v_0\}$ belong to $E(C)$. Equalities \eqref{trean1} and \eqref{s1} imply that
\[
w(\{u_1, v_0\}) = l(v_0) = w(\{u_2, v_0\}).
\]
Consequently \eqref{uro1} holds with $e_1 = \{u_1, v_0\}$, $e_2 = \{u_2, v_0\}$.
Thus $w : E(G) \to \mathbb{R}^+$ is pseudoultrametrizable by Theorem \ref{kasya}.
Theorem \ref{kasya} implies also that the mapping $\rho_w : V(G) \times V(G) \to \mathbb{R}^+$ 
defined by formula \eqref{kjvdde} is a pseudoultrametric on $V(G)$. 
Equality \eqref{zero} implies the equality $d_l=\rho_w$. Thus $d_l$ also is a pseudoultrametric on $V(G)$.



 Let $G$ be locally finite.
 Suppose that $d_l$ is an ultrametric on $V(G)$.
Let $\{u,v\}$ be an edge of $G$, 
\begin{equation*}
\{u,v\} \in E(G).
\end{equation*}
Then inequality \eqref{eq:10}  follows from Corollary \ref{btoy} and the definition of ultrametrics.

Suppose now that inequality \eqref{eq:10} holds for all $\{u,v\} \in E(G).$
We must show that $d_l$ is an ultrametric on $V(G).$

Suppose contrary that $d_l$ is not an ultrametric.
Then, using the definitions of the ultrametric spaces and pseudoultrametric ones, we can find different $u^*, v^* \in V(G)$ such that
\begin{equation}
\label{s4}
d_l(u^*,v^*)=0.
\end{equation}


If the inequality
\[
\max\{l(u^*),l(v^*)\}>0
\]
holds, then using \eqref{fdes1} with \( x=u^*\) and \( y=v^* \), we obtain
\[
d_l(u^*,v^*)>0,
\]
contrary to \eqref{s4}. Thus the equalities 
\begin{equation}
    \label{s5}
l(u^*)=0=l(v^*)
\end{equation}
hold.

Let \( N(u^*) \) be the set of all vertices of \( G \) which are adjacent with \( u^* \),
\[
N(u^*)=\{u\in V(G): \{u,u^*\}\in E(G)\}.
\]
Recall that the degree \( \delta_G(u^*) \) is the cardinal number of the set \( N(u^*) \),
\[
\delta_G(u^*)=|N(u^*)|.
\]
The graph \( G \) is locally finite. Consequently, the degree $\delta_G(u^*)$ is finite. Hence the set \( N(u^*) \) is finite,
\begin{equation}
    \label{s6}
|N(u^*)|<\infty.
\end{equation}
Moreover, since \eqref{eq:10} holds for all \( \{u,v\}\in E(G) \), double equality \eqref{s5} implies the inequality 
\[ 
l(u)>0 
\] 
for each \( u\in N(u^*) \).
The last inequality and \eqref{s6} give us the inequality
\begin{equation}
    \label{s7}
\inf_{u\in N(u^*)} l(u) > 0.
\end{equation}
Now using \eqref{s7} and formula \eqref{fdes1} we can prove that
\[
d_l(u^*,v^*) \ge \inf_{u\in N(u^*)} l(u) > 0,
\]
contrary to \eqref{s4}. 

Thus, if $G$ is locally finite and \eqref{eq:10} holds for each $\{u,v\}\in E(G),$ then pseudoultrametric \( d_l: V(G)\times V(G)\to \mathbb{R}^+ \) is an ultrametric on \( V(G) \) as required.

The proof is completed.
\end{proof}

\end{document}
